{\actuality} Декодирование внешнего стимула по~записям активности мозга "--- ключевая задача при~разработке интерфейсов мозг"--~компьютер и~изучении механизмов обработки визуальной информации. Регистрируемые нейросигналы представляют собой \textbf{мультимодальные многомерные временные ряды}: модальности имеют различные частоты дискретизации, взаимные задержки и~уровни шума, а~объем размеченной выборки ограничен длительностью сессии одного испытуемого. Функциональная магнитно-резонансная томография (фМРТ) обладает высоким пространственным разрешением при~низком временном, электроэнцефалография (ЭЭГ) "--- наоборот, поэтому совместное использование модальностей потенциально компенсирует ограничения каждой из~них.

Современные методы декодирования опираются на~контрастивные представления и~диффузионные генеративные модели, однако существующие подходы преимущественно ограничены одной модальностью и~требуют больших объемов обучающих данных. Модели глубокого обучения в~условиях малых выборок теряют обобщающую способность, что определяет необходимость методов, явно описывающих \textbf{пространственную и~временную структуру сигнала}: гемодинамический лаг отклика, локализацию информативных областей мозга и~геометрию ковариационных характеристик временных рядов.

\ifsynopsis\textit{Таким образом, актуальной задачей является разработка методов декодирования мультимодальных многомерных временных рядов, явно учитывающих их пространственно-временную структуру и~сохраняющих работоспособность в~условиях ограниченного объема выборки.}\else{Таким образом, актуальной задачей является разработка методов декодирования мультимодальных многомерных временных рядов, явно учитывающих их пространственно-временную структуру и~сохраняющих работоспособность в~условиях ограниченного объема выборки.}\fi

{\aim} данной работы является разработка методов извлечения пространственно-временных характеристик для~декодирования мультимодальных многомерных временных рядов, устойчивых в~условиях ограниченного объема выборки. Для достижения цели были поставлены и~решены следующие {\tasks}:

\begin{enumerate}[beginpenalty=10000]
    \item Формализовать задачу декодирования мультимодальных многомерных временных рядов с~учетом межмодальных задержек и~разнородности частот дискретизации.
    \item Разработать метод оценки гемодинамического лага BOLD-сигнала на~основе линейной авторегрессионной модели прогноза фМРТ-отклика по~визуальному стимулу.
    \item Разработать метод классификации временных рядов фМРТ на~основе масок активности и~римановой геометрии ковариационных матриц, устойчивый при~ограниченном объеме выборки одного испытуемого.
    \item Предложить мультимодальную архитектуру реконструкции визуальных стимулов по~одновременным сигналам фМРТ и~ЭЭГ на~основе контрастивного выравнивания совместных представлений с~пространством эмбеддингов изображений.
    \item Провести экспериментальную проверку предложенных методов на~открытых наборах данных и~исследовать границы их применимости при~доменных сдвигах.
\end{enumerate}

{\object} настоящей работы являются мультимодальные многомерные временные ряды нейросигналов (фМРТ, ЭЭГ), регистрируемые при~предъявлении визуальных стимулов.

{\subject} данной работы являются пространственно-временные характеристики нейросигналов: гемодинамический лаг отклика, маски активности областей мозга, ковариационные представления временных рядов на~римановом многообразии и~совместные мультимодальные представления, выравниваемые с~пространством эмбеддингов изображений.

{\novelty}
\begin{enumerate}[beginpenalty=10000]
    \item Задача декодирования визуальной информации по~сигналам мозга формализована в~единой постановке, объединяющей три уровня описания структуры данных: временную динамику, совместную пространственно-временную структуру и~мультимодальное объединение.
    \item Предложена линейная авторегрессионная модель прогнозирования фМРТ-отклика по~визуальному стимулу, обеспечивающая оценку гемодинамического лага BOLD-сигнала; полученная оценка согласуется с~известным физиологическим диапазоном.
    \item Разработан метод пространственно-временной классификации временных рядов фМРТ, основанный на~масках активности классов и~проекции ковариационных матриц на~касательное пространство риманова многообразия, размерность признакового описания которого не~зависит от~объема выборки.
    \item Впервые предложена мультимодальная архитектура реконструкции визуальных стимулов по~одновременным сигналам фМРТ и~ЭЭГ, объединяющая контрастивное выравнивание совместных представлений с~CLIP-пространством и~двухстадийную диффузионную генерацию.
\end{enumerate}

{\tinfluence} полученных результатов состоит в~том, что предложена единая постановка задачи декодирования мультимодальных многомерных временных рядов, в~рамках которой формализованы оценка межмодальной задержки, извлечение пространственно-временных признаков на~римановом многообразии и~контрастивное выравнивание разнородных модальностей; для~метода классификации показана независимость размерности признакового описания от~объема выборки.

{\pinfluence} настоящей работы заключается в~возможности применения разработанных методов в~неинвазивных интерфейсах мозг"--~компьютер, системах нейрореабилитации и~ассистивных технологиях. Устойчивость методов к~малым выборкам снижает требования к~объему дорогостоящих размеченных нейроданных. Программные реализации всех методов опубликованы в~открытом доступе и~обеспечивают воспроизводимость результатов.

{\methods} Для достижения поставленных
в~диссертационной работе целей используются:
\begin{enumerate}[beginpenalty=10000]
    \item Методы линейной алгебры и~математического анализа: линейные модели с~$L_2$-регуляризацией, спектральные разложения, риманова геометрия многообразия симметричных положительно определенных матриц.
    \item Методы теории вероятностей и~математической статистики: кросс-корреляционный анализ, проверка статистических гипотез с~поправками на~множественные сравнения.
    \item Методы машинного и~глубокого обучения: контрастивное обучение представлений, диффузионные генеративные модели, сверточные и~полносвязные архитектуры энкодеров.
    \item Методы вычислительного эксперимента: анализ исключением компонент, кросс-доменное тестирование, воспроизводимая постановка экспериментов.
\end{enumerate}

{\defpositions}
\begin{enumerate}[beginpenalty=10000]
    \item Единая постановка задачи декодирования мультимодальных многомерных временных рядов, объединяющая оценку межмодальной задержки, пространственно-временную классификацию и~мультимодальную реконструкцию.
    \item Линейная авторегрессионная модель прогнозирования фМРТ-отклика по~визуальному стимулу и~метод оценки гемодинамического лага BOLD-сигнала на~ее основе.
    \item Метод классификации временных рядов фМРТ на~основе масок активности классов и~проекции ковариационных матриц на~касательное пространство риманова многообразия, устойчивый в~условиях ограниченного объема выборки.
    \item Мультимодальная архитектура реконструкции визуальных стимулов по~одновременным сигналам фМРТ и~ЭЭГ с~контрастивным выравниванием совместных представлений и~двухстадийной диффузионной генерацией.
\end{enumerate}

{\passport} Диссертационная работа соответствует паспорту научной специальности 1.2.1 "--- <<Искусственный интеллект и машинное обучение>>, а именно в наибольшей степени следующим направлениям исследований: п.~2 <<Исследования в области оценки качества и эффективности алгоритмических решений для систем искусственного интеллекта и машинного обучения>>; п.~5 <<Методы и технологии обработки информации и извлечения знаний из различных источников>>; п.~15 <<Математические исследования в области статистики, логики, алгебры, топологии, анализа функций и других областях, ориентированные на решение задач искусственного интеллекта и машинного обучения>>; п.~16 <<Исследования в области специальных методов оптимизации, проблем сложности и снижения размерности>>; п.~17 <<Исследования в области многослойных алгоритмических конструкций, в том числе многослойных нейросетей>>.

{\reliability} полученных результатов обеспечивается корректностью
применения апробированного в научной практике математического аппарата линейной алгебры, теории вероятностей, математической статистики и римановой геометрии, а также экспериментальной проверкой разработанных методов на открытых наборах данных с проверкой статистической значимости результатов. Для проведенных вычислительных экспериментов даются детальные описания, обеспечивающие их воспроизводимость; программные реализации опубликованы в открытом доступе. Результаты опубликованы в рецензируемых научных изданиях и в трудах рецензируемых российских конференций по машинному обучению и анализу данных.

{\probation} Основные результаты работы докладывались и обсуждались на следующих научных конференциях:
\begin{enumerate}
    %% \item Международная конференция <<Интеллектуализация обработки информации (ИОИ-2026)>>, сентябрь 2026~г.
    \item 68-я Всероссийская научная конференция МФТИ,
    Долгопрудный, \fixme{DD--DD месяц} 2026~г.
    \item Конференция AI Journey, Москва, 19--21 ноября 2025~г.
    \item 22-я Всероссийская конференция с международным участием
    <<Математические методы распознавания образов (ММРО-2025)>>,
    Муром, 22--26 сентября 2025~г.
    \item 67-я Всероссийская научная конференция МФТИ,
    Долгопрудный, 31 марта -- 5 апреля 2025~г.
    \item 66-я Всероссийская научная конференция МФТИ,
    Долгопрудный, 1--6 апреля 2024~г.
\end{enumerate}


{\contribution} Все приведенные результаты, кроме отдельно оговоренных случаев, получены диссертантом лично при научном руководстве к.ф.-м.н.~А.\,В.~Грабового.

\ifnumequal{\value{bibliosel}}{0}
{%%% Встроенная реализация с загрузкой файла через движок bibtex8. (При желании, внутри можно использовать обычные ссылки, наподобие `\cite{vakbib1,vakbib2}`).
    {\publications} Основные результаты по теме диссертации изложены
    в~XX~печатных изданиях,
    X из которых изданы в журналах, рекомендованных ВАК,
    X "--- в тезисах докладов.
}%
{%%% Реализация пакетом biblatex через движок biber
    \begin{refsection}[bl-author, bl-registered]
        % Это refsection=1.
        % Процитированные здесь работы:
        %  * подсчитываются, для автоматического составления фразы "Основные результаты ..."
        %  * попадают в авторскую библиографию, при usefootcite==0 и стиле `\insertbiblioauthor` или `\insertbiblioauthorgrouped`
        %  * нумеруются там в зависимости от порядка команд `\printbibliography` в этом разделе.
        %  * при использовании `\insertbiblioauthorgrouped`, порядок команд `\printbibliography` в нем должен быть тем же (см. biblio/biblatex.tex)
        %
        % Невидимый библиографический список для подсчета количества публикаций:
        \phantom{\printbibliography[heading=nobibheading, section=1, env=countauthorvak,          keyword=biblioauthorvak]%
        \printbibliography[heading=nobibheading, section=1, env=countauthorwos,          keyword=biblioauthorwos]%
        \printbibliography[heading=nobibheading, section=1, env=countauthorscopus,       keyword=biblioauthorscopus]%
        \printbibliography[heading=nobibheading, section=1, env=countauthorconf,         keyword=biblioauthorconf]%
        \printbibliography[heading=nobibheading, section=1, env=countauthorother,        keyword=biblioauthorother]%
        \printbibliography[heading=nobibheading, section=1, env=countauthor,             keyword=biblioauthor]%
        \printbibliography[heading=nobibheading, section=1, env=countauthorvakscopuswos, filter=vakscopuswos]%
        \printbibliography[heading=nobibheading, section=1, env=countauthorscopuswos,    filter=scopuswos]}%
        %
        \nocite{*}%
        %
        {\publications} \ifsynopsis\textit{Список публикаций приведен в конце автореферата. }\else\textit{Список публикаций приведен в конце диссертации. }\fi Основные результаты по теме диссертации изложены в~\arabic{citeauthor}~научных работах, из которых~\arabic{citeauthorscopuswos} "--- в~периодических научных журналах, индексируемых Web of~Science и Scopus, а остальные~\arabic{citeauthorconf} "--- в~тезисах докладов российских и международных конференций.
    \end{refsection}%
    \begin{refsection}[bl-author, bl-registered]
        % Это refsection=2.
        % Процитированные здесь работы:
        %  * попадают в авторскую библиографию, при usefootcite==0 и стиле `\insertbiblioauthorimportant`.
        %  * ни на что не влияют в противном случае
        \nocite{dorin2024forecasting}
        \nocite{dorin2025enhancing}
        \nocite{dorin2026decoding}
        \nocite{dorin2025mmro}
        \nocite{dorin2024conf66mipt}
        \nocite{dorin2025conf67mipt}
        \nocite{dorin2026conf68mipt}
        \nocite{dorin2025aij}
    \end{refsection}%
}
